\documentclass[11pt,a4paper]{article}

\usepackage{../shared/cathedral-preamble}

\title{%
  Open Horizons:\\
  Research Frontiers Beyond the Cathedral%
}
\author{Jason Robert Gochanour}
\date{April 28, 2026}

\begin{document}
\maketitle

\begin{abstract}
We survey the open problems and research directions suggested by
the Cathedral project, spanning computer science, pure mathematics,
and formalization engineering. The Cathedral reduced the Riemann
Hypothesis to two axioms; we identify the mathematical infrastructure
needed to close them, the algorithmic questions raised by the
spectral pipeline, and the broader implications for the formal
verification of analytic number theory.
\end{abstract}

\tableofcontents

% ================================================================
\section{Introduction}
% ================================================================

The Cathedral is complete as a proof architecture, but it opens
more doors than it closes. This paper maps the territory beyond
the Cathedral walls.

% ================================================================
\section{Closing the Two Axioms}
% ================================================================

\subsection{Axiom 1: Hardy--Littlewood MVT}

Formalizing $\int_0^T|1/\zeta(1/2+it)|^2\,dt = O(T)$ requires:
\begin{enumerate}
  \item The approximate functional equation for $\zeta(s)$.
  \item Fourth-moment bounds on the critical line.
  \item Hybrid Euler--Hadamard product estimates.
\end{enumerate}

\textbf{Difficulty}: Hard. Requires 3--6 months of Mathlib
development by an expert in analytic number theory formalization.

\subsection{Axiom 2: Hadamard Product Lower Bound}

Formalizing $|\zeta(s)| \geq c|t|^{-A}$ for $\re s \geq 1/2 + \varepsilon$ requires:
\begin{enumerate}
  \item The Hadamard product formula for $\zeta(s)$.
  \item Riemann--von Mangoldt zero counting $N(T) \sim T\log T/(2\pi)$.
  \item Connecting Borel--Carath\'eodory (in Mathlib) to zero density.
\end{enumerate}

\textbf{Difficulty}: Medium--Hard. Borel--Carath\'eodory is in
Mathlib; the gap is infrastructure.

% ================================================================
\section{Algorithmic Questions}
% ================================================================

\subsection{Fast Gram Matrix Computation}

Computing $G_N$ exactly via Vasyunin's formula costs $O(N^2 \log N)$.
Can this be improved? The Fourier structure of the cotangent sums
suggests an $O(N \log^2 N)$ algorithm via FFT.

\subsection{The Eigenvalue Distribution}

Numerical experiments show $\lambda_{\min}(G_N) \sim N^{-2}$ while
$\lambda_{\max}(G_N) \sim \log N$. Can these asymptotics be proved?
The connection to random matrix theory (GOE spacing statistics)
suggests tools from integrable probability.

\subsection{Optimal Weight Computation}

The optimal BD weights $v = G^{-1}b$ grow like $\|v\|^2 = \Theta(N)$
while $d_N^2 \to 0$. The Sherman--Morrison rank-1 update
$G_{N+1} = G_N + uu^T$ enables $O(N^2)$ online updating. Can this
be reduced further?

% ================================================================
\section{Mathematical Frontiers}
% ================================================================

\subsection{The Spectral Universality Conjecture}

Exploration 19 established experimentally that the spectral
statistics of the BD Gram matrix transition from Poisson (small $N$)
to GOE (large $N$) at a critical threshold
$N_c \approx 60 \cdot m/\varphi(m)$.

\textbf{Open problem}: Prove the universality of this threshold.
This would connect the Cathedral to the Bohigas--Giannoni--Schmit
conjecture in quantum chaos.

\subsection{The $\alpha \approx 0.47$ Constant}

The participation ratio stabilizes at $\alpha \approx 0.47$,
between full localization and GOE delocalization.

\textbf{Open problem}: Derive $\alpha$ from the zeta function.
This may involve the pair correlation of zeta zeros
(Montgomery's conjecture).

\subsection{Generalized Nyman--Beurling}

The Cathedral uses $h_k(x) = \{1/(kx)\}$. What happens with:
\begin{itemize}
  \item $h_k(x) = \{k^\alpha/x\}$ for $\alpha \neq 1$?
  \item $h_k(x) = \{1/(kx)\} \cdot w(x)$ for weights $w$?
  \item Analogues for Dirichlet $L$-functions?
\end{itemize}

% ================================================================
\section{Formalization Engineering Frontiers}
% ================================================================

\subsection{Mathlib Contributions}

Five Mathlib additions with highest impact:
\begin{enumerate}
  \item Vertical contour integrals.
  \item Analytic continuation for Mellin transforms.
  \item The Mertens function with basic bounds.
  \item Fractional part API enrichment.
  \item Montgomery--Vaughan mean value theorem.
\end{enumerate}

\subsection{Certified Computation at Scale}

The cert-export pipeline demonstrates JSON-to-Lean certification.
Scaling to $N = 10{,}000$ requires:
\begin{itemize}
  \item GPU-accelerated MPFR eigensolvers.
  \item Parallel certification across moduli.
  \item Streaming certificate generation for Lean.
\end{itemize}

\subsection{Formal Proofs for Related Problems}

The Cathedral architecture could be adapted for:
\begin{itemize}
  \item The Generalized Riemann Hypothesis (Dirichlet $L$-functions).
  \item The Selberg class (axiomatic $L$-functions).
  \item Robin's inequality---already proved equivalent to NB
    in the Cathedral (\S5.1 of the main paper). Extending to
    a standalone formalization independent of the Crown axioms
    would be of interest.
\end{itemize}

% ================================================================
\section{Conclusion}
% ================================================================

The Cathedral is an ending and a beginning. The proof architecture
is complete; the mathematical questions it raises are not. The
research frontier extends in at least five directions: closing the
two axioms, understanding the spectral universality, exploiting the
algorithmic structure, contributing to Mathlib, and generalizing to
other $L$-functions.

\begin{thebibliography}{99}
\bibitem{cathedral}
J.R.~Gochanour, \emph{The Cathedral}, 2026.
\bibitem{montgomery1973}
H.L.~Montgomery, \emph{The pair correlation of zeros of the zeta
function}, Proc. Symp. Pure Math. \textbf{24} (1973).
\bibitem{bgs1984}
O.~Bohigas, M.J.~Giannoni, and C.~Schmit, \emph{Characterization
of chaotic quantum spectra}, Phys. Rev. Lett. \textbf{52} (1984).
\end{thebibliography}

\end{document}
